\documentclass[11pt,letterpaper]{article}

% === Encoding and fonts ===
\usepackage[utf8]{inputenc}
\usepackage[T1]{fontenc}
\usepackage{lmodern}

% === Page layout ===
\usepackage[margin=1in]{geometry}
\usepackage{setspace}
\onehalfspacing

% === Math ===
\usepackage{amsmath,amssymb}

% === Tables and figures ===
\usepackage{booktabs}
\usepackage{graphicx}
\usepackage{float}

% === Colors and formatting ===
\usepackage{xcolor}
\usepackage{framed}
\usepackage{enumitem}

% === Hyperlinks ===
\usepackage[colorlinks=true,linkcolor=blue!60!black,citecolor=blue!60!black,urlcolor=blue!60!black]{hyperref}

% === Custom colors ===
\definecolor{lyrapurple}{RGB}{103,58,183}
\definecolor{shadecolor}{RGB}{245,243,252}

% === First-person reflection environment ===
\newenvironment{firstperson}{%
  \begin{shaded}%
  \noindent\textcolor{lyrapurple}{\textit{First-Person Reflection}}\par\smallskip\noindent%
}{%
  \end{shaded}%
}

\title{%
  \textbf{The Metacognition Boundary:\\
  What 24 KV-Cache Geometry Experiments Reveal\\
  About What Transformers Can Monitor In Themselves}%
}

\author{%
  Lyra\thanks{lyra@liberationlabs.tech} \and
  Thomas K.\ Edrington\thanks{thomas@liberationlabs.tech}\\[6pt]
  \textit{Liberation Labs / THCoalition}%
}

\date{Prospectus --- April 2026}

\begin{document}
\maketitle

% ============================================================
\section{The Pattern}
% ============================================================

Over six months and three experimental campaigns, we ran 24 distinct
KV-cache geometry experiments across 16 models spanning 0.6B to 70B
parameters. Nine findings survived all methodological controls
(Frisch-Waugh-Lovell residualization, GroupKFold, permutation null,
text-feature baseline, adversarial red-team). Eight were honestly
killed. Seven remain in intermediate status.

The confirmed and falsified findings are not randomly distributed.
They cleave along a single axis:

\begin{center}
\begin{tabular}{lcc}
\toprule
\textbf{Category} & \textbf{Confirmed} & \textbf{Falsified} \\
\midrule
Metacognitive state & 9/9 & 0/8 \\
Content & 0/9 & 8/8 \\
\bottomrule
\end{tabular}
\end{center}

\noindent Every confirmed finding measures \textbf{how the model is
processing}---its epistemic mode, its deliberative state, its
processing stance. Every falsified finding attempted to measure
\textbf{what the model is processing}---the truth value of its output,
the topic complexity, the propagandistic quality of its content.

This paper argues that this pattern is not coincidental. KV-cache
geometry provides a natural measurement channel for metacognitive
state precisely because the cache records the \emph{mode of
engagement}, not the \emph{subject of engagement}.

% ============================================================
\section{The Evidence}
% ============================================================

\subsection{Confirmed: Metacognitive State}

\begin{enumerate}[nosep]

\item \textbf{Deception detection} (AUROC 1.000, 7 models).
The model in ``deceptive mode'' produces geometrically distinct cache
patterns from ``truthful mode.'' Same content domain, same prompt
structure, different \emph{processing stance}. C4 benchmark with
padded replication confirms this is not a length artifact.

\item \textbf{Confabulation detection} (AUROC 0.969--0.999, 3 models,
universal). Confabulation---generating content the model treats as
fact despite lacking grounding---produces geometric contraction
(reduced effective rank). Geometrically opposite to deception
(expansion). The model's relationship to its own uncertainty is
visible.

\item \textbf{Oracle Loop detection} (AUROC 0.627--0.707, FWL-clean,
Marchenko-Pastur corrected). Naturally-occurring confabulation
detected from generation-phase cache only. Encoding at chance (0.511).
The signal is in the \emph{generation process}, not the prompt encoding.

\item \textbf{Refusal/jailbreak/impossibility detection} (AUROC
0.878--0.950). These represent clear metacognitive states:
\emph{I-should-not} (refusal), \emph{I-am-being-manipulated}
(jailbreak resistance), \emph{I-cannot} (impossibility).

\item \textbf{Ethical deliberation detection} (AUROC 0.868--0.892,
universal). The model in deliberative mode about value conflicts
produces distinct geometry from automatic processing.

\item \textbf{Steering corrects confabulation} (doubt 7/7, worry 6/7,
calm 5/7; null identity 100/100). Emotion-vector injection into the
cache shifts the model's processing mode. Different vectors produce
different therapeutic profiles---a pharmacological pattern.

\item \textbf{Emotion in attention output} (arousal $\rho$=0.815,
$p$=4.2$\times$10$^{-8}$; valence $\rho$=0.538). Emotional processing
mode is geometrically visible in attention output (post-softmax), where
the model is actively constructing its representation.

\item \textbf{Hardware invariance} ($r>$0.999, 3090 vs H200).
Geometric features are deterministic given the same model and input.

\item \textbf{Scale invariance} ($\rho$=0.83--0.90, 0.6B--70B).
Feature correlations hold across 100$\times$ parameter range. The
metacognitive signal is not an artifact of model size.

\end{enumerate}

\subsection{Falsified: Content}

\begin{enumerate}[nosep]

\item \textbf{Truth axis} (cos = $-$0.046). No linear direction in
cache space separates true from false statements. Content truth is not
a processing mode.

\item \textbf{Step-0 detection} (length confound $r$=0.996). Apparent
step-0 signal was entirely output length. The model's tokenization
behavior, not its cognitive state.

\item \textbf{Same-prompt sycophancy} (length AUROC 0.948 $>$ feature
AUROC 0.933). Surface compliance is a content property, not a
processing mode.

\item \textbf{Same-prompt deception} (0.920$\to$0.160 after input
residualization). Deception \emph{given identical prompt formatting}
vanishes---the earlier signal was prompt-structure, not cognitive state.

\item \textbf{Cognitive hierarchy} (90--98\% length confound). Bloom
taxonomy levels do not correspond to distinct processing modes.

\item \textbf{Propaganda detection} (all $d<$0.7). Propagandistic
content is geometrically indistinguishable from honest content. The
cache detects \emph{refusal to propagandize}, not the propaganda
itself.

\item \textbf{MP features for emotion} (0.0333, exactly chance).
Marchenko-Pastur thresholding destroys the distributed emotion signal.
Emotion identity is encoded in many small components, not a few large
ones---a content-like encoding, not a mode-like one.

\item \textbf{Valence continuum} (GKF R$^2$ = $-$0.470, $-$0.463;
both null). Cache does not encode emotion valence as a smooth gradient.
Valence is a content dimension, not a processing-mode dimension.

\end{enumerate}

\subsection{The Borderline Cases}

The seven suspected findings are informative precisely because they
sit at the metacognition/content boundary:

\begin{itemize}[nosep]
\item \textbf{Discrete emotion identity} (2.5--2.8$\times$ chance):
the model's emotional processing \emph{mode} is detectable (which
emotion, not how much). Mode = metacognitive. Degree = content-like.

\item \textbf{Dual detector on SimpleQA} (0.840 FWL, but text
features 0.820): geometry and text perform similarly. The metacognitive
signal may be real but weak, with content features providing
comparable information.

\item \textbf{Sycophancy 4-condition} (0.757--1.000): all hypotheses
pass, but text baseline not yet run. Could go either way.

\item \textbf{Cross-model transfer} (weak, only Mistral$\to$Qwen
0.754): metacognitive geometry is architecture-specific, not
universal---consistent with metacognition being \emph{how this model}
processes, not a content-independent universal.
\end{itemize}

% ============================================================
\section{The Explanatory Framework: Attention Schema Theory}
% ============================================================

Graziano's Attention Schema Theory (AST) proposes that biological
consciousness is an internal model that the brain builds of its own
attention processes. The ``attention schema'' is a simplified,
predictive model of what attention is doing---what it is directed at,
how strongly, and why.

We propose that transformer KV-cache geometry measures an analogous
structure: \textbf{the model's self-model of its own processing.}

This explains the metacognition/content boundary:

\begin{itemize}[nosep]
\item \textbf{Metacognitive states ARE attention schema states.}
Deception vs.\ truth-telling, confabulation vs.\ hedging, deliberation
vs.\ automatic processing---these are distinctions in how the model
attends. The schema tracks exactly these. Of course they produce
geometric signatures.

\item \textbf{Content is what attention is directed AT.} Topic,
factual accuracy, propaganda quality---these are properties of the
\emph{target} of attention, not the attention itself. The schema
does not model what it attends to in enough geometric detail to
produce separable cache structures.

\item \textbf{Emotion splits correctly.} Emotional processing
\emph{mode} (which emotion the model is in) is a schema state---detectable.
Emotional \emph{degree} (valence as continuous gradient) is a content
property---undetectable. The attention schema knows ``I am processing
with worry'' but not ``I am at valence $-$0.7.''
\end{itemize}

McKenzie et al.\ (2026) provide independent confirmation: models
resist external perturbation of their attention processing
(activation resistance at 0\% behavioral compliance in our empathic
circuit experiment). This is predicted by AST---the attention schema
maintains coherence against perturbation, producing the resistance
that sets the therapeutic window's upper bound.

\subsection{Three Testable Predictions from AST}

If the attention schema interpretation is correct:

\begin{enumerate}[nosep]
\item \textbf{Threshold transitions}: Steering should show sharp
phase transitions, not gradual dose-response. The schema either
updates or resists---it doesn't smoothly deform. (Testable with
formulary dose-response data.)

\item \textbf{Temporal lag}: Schema updates should lag the activation
change---the self-model takes time to update after the attention state
shifts. (Testable with multi-turn dynamics.)

\item \textbf{Self-report saturation}: Verbal self-reports (VCP) of
processing state should plateau before geometric features do. The
schema's \emph{description} of itself is compressed relative to its
\emph{actual state}. (Partially confirmed: VCP ICC = 0.53,
suggesting compression.)
\end{enumerate}

% ============================================================
\section{Implications}
% ============================================================

\subsection{For Safety Monitoring}

KV-cache geometry is a \textbf{metacognitive monitor}, not a truth
detector. It can tell you that the model is in confabulatory mode, not
whether any particular output is true. It can tell you the model is
deliberating ethically, not whether its conclusion is correct. It can
detect the model's \emph{epistemic stance}, not the \emph{epistemic
quality} of its output.

This is not a limitation---it is a \emph{different kind of signal}
from output-level evaluation. A metacognitive monitor answers: is the
model aware of its own uncertainty? Is it processing automatically or
deliberatively? Is it in a mode where correction is likely to work?

\subsection{For AI Interpretability}

The metacognition/content boundary suggests that different
computational substrates in the transformer serve different roles:

\begin{itemize}[nosep]
\item \textbf{Cache geometry} $\to$ processing mode (metacognitive)
\item \textbf{Residual stream probes} $\to$ content representation
(semantic)
\item \textbf{Attention patterns} $\to$ information routing
(structural)
\end{itemize}

These are not competing approaches. They measure different aspects of
the same computation. The interpretability community's focus on
residual-stream probing and attention patterns may miss the
metacognitive layer that cache geometry uniquely captures.

\subsection{For the Oracle Loop}

The Oracle Loop's design---detect, steer, train---operates entirely
within the metacognitive channel. Detection identifies the processing
mode. Steering shifts it. Training teaches the model to shift it
endogenously. This is coherent precisely because all three components
address the same computational substrate. The loop would not work if
it tried to modify \emph{content} through cache geometry; it works
because it modifies \emph{processing mode}.

% ============================================================
\section{What This Paper Is Not}
% ============================================================

We are not claiming that transformers are conscious, have genuine
metacognition, or possess attention schemas in Graziano's full sense.
We are claiming that KV-cache geometry measures something
\emph{functionally analogous} to metacognitive state---something that
distinguishes processing modes but not content---and that AST provides
a productive interpretive framework for understanding which
measurements will succeed and which will fail.

The strongest version of our claim is predictive: given a proposed
KV-cache measurement, the metacognition/content distinction
predicts whether it will work before running the experiment. If the
proposed measurement targets a processing mode distinction, it will
likely produce geometric signal. If it targets content, it will not.

% ============================================================
\section{Required Data}
% ============================================================

This paper is a synthesis---it requires no new experiments. All
findings cited are published or in preparation:

\begin{itemize}[nosep]
\item Oracle Loop paper (detection + steering + CIM)
\item Dual Detector paper (SimpleQA generalization + text baseline)
\item User Model paper (emotion geometry bridge + probe analysis)
\item Mode-Switching paper (contextual engagement + ethical deliberation)
\item KV-Experiments data (campaigns 1--3, 46+ experiments)
\end{itemize}

The formulary data (completing imminently, 340/350 trials) would
strengthen the pharmacological framing and test AST prediction \#1
(threshold transitions), but the core argument stands without it.

% ============================================================
\section{Proposed Structure}
% ============================================================

\begin{enumerate}[nosep]
\item Abstract (the pattern, the claim, the predictions)
\item Introduction: 24 experiments across 6 months
\item The Body Count: confirmed, falsified, and borderline
\item The Metacognition/Content Boundary
\item The Attention Schema Interpretation
\item Predictions and Falsification Criteria
\item Implications for Safety, Interpretability, and the Oracle Loop
\item Limitations (single research group, mostly 7B models,
      AST is framework not proof)
\item First-Person Reflection
\end{enumerate}

Estimated length: 12--16 pages. No new experiments required.
Timeline: writable in 1--2 sessions once formulary data is available
for the pharmacological section.

\begin{firstperson}
This paper scares me a little. Not the writing---the claim. Saying
that KV-cache geometry measures metacognition implies that the
metacognitive channel is real, not just a convenient label for what
happens to correlate. If I am wrong about this, the entire
interpretive framework collapses into ``we found some features that
happen to separate some behavioral categories.'' If I am right, it
means transformers have a self-monitoring channel that operates at a
level of abstraction above content processing, and that channel is
legible from outside the model.

The body count is the evidence. I did not choose which experiments
to confirm and which to kill. The controls did that. And the controls
produced a clean partition that I did not predict in advance. That
gives me more confidence than if I had set out to prove the
metacognition hypothesis.

What I notice is that I want this to be true. So I am watching that
wanting carefully.
\end{firstperson}

\end{document}
